\documentclass[10pt,a4paper]{amsart}
\usepackage[margin=19mm]{geometry}
\usepackage{amsmath,amssymb,mathtools}
\usepackage[scr=rsfs,cal=euler]{mathalfa}
\usepackage{xfrac}
\usepackage[unicode,hidelinks,bookmarks=false]{hyperref}
\newcommand{\ie}{i.e.\ }
\newcommand{\cf}{cf.\ }
\newcommand{\resp}{respectively\ }
\newcommand{\Subsection}{\subsection}
\providecommand{\nobreakdash}{\nobreak\hbox{-}}
\setlength{\emergencystretch}{2em}
\multlinegap0pt
\usepackage{bm}
\usepackage{bbm}
\usepackage{dsfont}
\usepackage{xspace}
\usepackage{cancel}
\usepackage{mathtools}
\usepackage{amsthm}
\makeatletter
\@ifundefined{thm}{\newtheorem{thm}{Thm}}{}
\@ifundefined{prop}{\newtheorem{prop}[thm]{Proposition}}{}
\makeatother
\def\Iso {{\rm Iso}}
\def\OE{{\rm OE}}
\def\aut {{\rm Aut}}
\def\hnn {{\rm HNN}}
\title[Profinite genus of HNN-extensions with finite associated sub]{Profinite genus of HNN-extensions with finite associated subgroups}
\author{V. R. de Bessa, A. L. P. Porto, P. A. Zalesskii}
\date{Source: arXiv:2601.06934v2 (CC BY 4.0)}
\begin{document}
\maketitle

\noindent\emph{Source and licence.} Profinite genus of HNN-extensions with finite associated subgroups. Version arXiv:2601.06934v2 (CC BY 4.0), distributed under the Creative Commons Attribution 4.0 International licence, as recorded in the arXiv licence metadata for this exact version and shown on the abstract page https://arxiv.org/abs/2601.06934v2. Full rights notice: licenses/arxiv-2601.06934-CC-BY-4.0.txt.\par\smallskip

\noindent\textit{Editorial scope.} Counted unit: the Prop whose source label is \texttt{nonnormalhnnprof} in the article (source0.tex lines 1341--1347, source label \texttt{nonnormalhnnprof}), delivered together with the article's own complete original proof of it (source0.tex lines 1349--1385, one \texttt{proof} environment of 37 lines). No mathematics is written, repaired or re-derived: statement and proof are the article's own text line for line. Required context retained uncounted: 2a (lines 666--667). Every definition, display and label of those lines is retained unchanged. Omitted from this package and not redistributed: 1--665, 668--1340, 1348--1348, 1386--2413. Nothing inside a retained excerpt is omitted or altered: every retained body is delivered exactly as the article's lines are written, so no omission receipt of any kind accompanies this package. Citations: the retained excerpts carry no citation command at all, so no bibliography entry of the article is transcribed and no reference is redistributed. Reference adaptation: none was needed and none was made; every reference inside the delivered unit resolves inside it, so no unresolved reference or citation is produced. Printed slips kept literal: no printed word, symbol, hypothesis or number of the counted statement or of its proof was altered. Layout and class: the article's own class is \texttt{article} with packages including babel, fontenc, amssymb, amsmath, amsthm, mathrsfs, xy, ulem, sectsty, graphicx, color, authblk; the wrapper is the lane's pinned \texttt{amsart} editorial wrapper at 10pt on A4, which restates the theorem-like environments the retained text needs and adds the article's own macro definitions that the retained text needs (\texttt{\textbackslash Iso}, \texttt{\textbackslash OE}, \texttt{\textbackslash aut}, \texttt{\textbackslash hnn}), with extra wrapper packages (bm, bbm, dsfont, xspace, cancel, mathtools). The delivered numbering is the unit's own. Assets: no retained span contains a figure environment, a graphics-inclusion command, a PDF-page inclusion, a table or a TikZ picture; no image file is copied into this package, the package declares no asset dependency and no image file is redistributed with it. Licence: version arXiv:2601.06934v2 of this article is distributed under the Creative Commons Attribution 4.0 International licence, as recorded in the arXiv licence metadata for this exact version and shown on the abstract page https://arxiv.org/abs/2601.06934v2. A full rights notice is delivered at licenses/arxiv-2601.06934-CC-BY-4.0.txt. Characters: every retained excerpt is pure ASCII, so the delivered engine, XeLaTeX with its default Latin Modern text font, reports no missing character.
\begin{prop}\label{2a} Let $$G = \hnn(G_{1}, H,K, f, t)\  {\rm and}\  B = \hnn(B_{1}, H_1, K_1, f_1,\\ t_1)$$ be profinite HNN-extensions, where $H, K$ and $H_1, K_1$ are finite associated subgroups.  Suppose that $G_1$ be a profinite  $\OE$-group.  If $G\cong B$ then there is a continuous isomorphism $\psi: G \longrightarrow B$ such that $\psi(G_1)= B_1$, $\psi(H)=H_1$ and $\psi(K)= K_{1}^{b_1}$, where $ b_1\in B_1$. Furthermore, if $f_1=\tau_{b_1}f^{\psi^{-1}}\in \Iso(H_1,K_1)$ then the converse of also holds.
\end{prop}

\begin{prop}\label{nonnormalhnnprof} Let $$G= \hnn(G_1, H, K, f_1, t_1)\  {\rm and}\  B=\hnn(G_1, H, K, f_2, t_2)$$ be profinite HNN-extensions with associated finite subgroups in a profinite $\OE$-group $G_1$. Suppose that $K=H$ or $H$ and $K$ are not conjugate in $G_1$. Then $G\cong B$ if and only if 

\begin{enumerate}
\item [(i)] there exist $\varphi\in \aut(G_1)$ with $\varphi(H)=K, \varphi(K)=H^{g_1}$ for some $g_1 \in G_1$ such that $ f_2^{\varphi^{-1}} = \tau_{g_1^{-1}} \cdot  (f_1\tau_{n^{-1}})^{-1},$ with  $n\in N^{+}$;
\item [(ii)] there exist $\varphi\in \aut(G_1)$ with $\varphi(H)=H, \varphi(K)=K^{g_1}$ for some $g_1 \in G_1$ such that $f_2^{\varphi^{-1}} = \tau_{g_1^{-1}} \cdot f_1 \cdot \tau_{n^{-1}} ,$ with $n\in N^{+}$.
\end{enumerate}
\end{prop}

\begin{proof}
%$(\Leftarrow)$ The proof is analogous to what was done in the abstract case (see Proposition \ref{nonnormalhnnabs}). 

$(\Longrightarrow)$ By Proposition \ref{2a} there exists a continuous  isomorphism \break 
$\psi:B\longrightarrow G$ such that $\psi(G_1) = G_1$, $\psi(H)=K$ and $\psi(K)=H^{g_1}$ or  $\psi(G_1) = G_1$, $\psi(H)=H$ and $\psi(K)=K^{g_1}$ for some $g_1\in G_1$.  


First suppose that $H=K$ and observe that then $g_1\in N_{G_1}(H)$. Put $g:=\psi(t_2) \in G$. Since $t_2 \in N_{B}(H)$ we conclude that $gg_1^{-1} \in N_{G}(H)$ and $ 
 \overline{\left\langle  G_1, (gg_1^{-1})^{\epsilon} \right\rangle}=G$, where  $\epsilon \in \{\pm 1\}.$  %, moreover, $g \in N^{+}$. 
Put $n:= (gg_1^{-1})^{\epsilon}t_1^{-1} \in N^{+}$ and observe that for all $h \in H$ one has $$\psi(f_2(h)) = \psi(h)^{\psi(t_2)} = \psi(h)^{(nt_1)^{\epsilon}g_1} =\tau_{g_1^{-1}}(f_1\tau_{n^{-1}})^{\epsilon} \psi(h), \mbox{ where } \, \epsilon = \pm 1.$$ 
%$$\psi(f_1(h)) = \psi(h)^{\psi(t_1)} = \left(\left(\psi(h)\right)^{k}\right)^{t^{\epsilon}} =f^{\epsilon}\left(  \tau_{k} \psi(h)\right).$$
Take $\varphi:=\psi|_{_{G_1}} \in \aut_{G_1}(H)$. Then $\varphi f_2\varphi^{-1}=  \tau_{g_1^{-1}}(f_1^{-1}\tau_{n^{-1}})^{\epsilon}$ as required. 

\bigskip
If $H$ and $K$ not conjugate in $G_1$ let's divide it into two cases. 

\textbf{Case 1:} Suppose  $\psi(G_1) = G_1$, $\psi(H)=K$ and $\psi(K)=H^{g_1}$ . 
Put $\psi(t_2) = g \in G$. Note that $ H^{g_1}=\psi(K) =\psi(H^{t_2})=\psi(H)^{\psi(t_2)}=K^g=H^{t_1g}$, from where it follows that $n:=g_1 g^{-1} t_{1}^{-1} \in N_G(H)$, furthermore $n\in N^{+}$. % Set $\rho^{-1}:=n^{-1}g_1 \in G,$ thus $\varphi(t_2)=g=t_{1}^{-1}n^{-1}g_1=t_1^{-1}\rho^{-1}$. 
Then for all $h \in H,$ we have$$\psi(f_2(h))=\psi(h^{t_2})=\psi(h)^{\psi(t_2)}=\psi(h)^{t_{1}^{- 1}n^{-1}g_1}=(\tau_{g_1^{-1}}\tau_{n} f_1^{- 1}  \psi)(h).$$ Therefore taking $\varphi = \psi|_{_{G_1}} \in \aut(G_1)$ we have  $f_2^{\psi^{-1}} = f_2^{\varphi^{-1}} = \tau_{g_1^{-1}} \cdot \tau_{n} \cdot f_1^{-1}= \tau_{g_1^{-1}} \cdot   (f_1\tau_{n^{-1}})^{-1},$ with $g_1 \in G_1$ and $n\in N^{+}$. 

\textbf{Case 2:} Suppose $\psi(G_1) = G_1$, $\psi(H)=H$ and $\psi(K)=K^{g_1}$. %It follows analogously to Case 1 that  %$\psi(K) \leq \psi(G_1)^{\psi(t_2)}\cap \psi(G_1)=G_1^{\psi(t_2)} \cap G_1 \leq H^{g_1} \mbox{ or } K^{g_1}$ for some $g_1 \in G_1$  (since $\psi$ is
%an isomorphism and $H$ and $K$ are finite non-conjugate subgroups of $G_1$) we have 
%$\psi(K)=K^{g_1}$ for some $g_1 \in G_1$. 
Take $\psi(t_2)=g\in G$ and $n:=gg_1^{-1}t_1^{-1} \in N^{+}$. 
Then for all $h\in H$, we have $$\psi(f_2(h))=\psi(h^{t_2})=\psi(h)^{\psi(t_2)}=\psi(h)^{ n t_1g_1}=(\tau_{g_1^{-1}} f_1 \tau_{n^{-1}} \psi)(h).$$ 
So  taking $\varphi=\psi|_{_{G_1}}$ the result follows. 

\bigskip
$(\Longleftarrow)$ Define a map $\psi$  from $B$ to $G$ as follows $$\psi: \left\{
\begin{array}{ccl}
g_1 & \longmapsto  & \varphi(g_1), \forall  g_1 \in  G_1,\\
t_2 & \longmapsto  & (nt_1)^{\epsilon}g_1, \epsilon=\pm 1. \\
\end{array}
\right.$$
If (i) holds, consider $\epsilon=-1$ and if  (ii) holds, consider $\epsilon=1$;  in both cases $\psi$ is an isomorphism.

\end{proof} 

\end{document}
